% Purpose: Methods 2.1, the data source, the target, the dataset, the descriptor classes and the domain statement (the single home of the domain limit, docs/story_spec.md (ii)). Input by sections/methods.tex.
\subsection{Data and Domain}
\label{sec:data}

All formation energies, descriptors and structures come from the database of neutral and charged oxygen vacancies of Kumagai et al.\ \citep{Kumagai2021insights}, released under the \ccby{} licence and pinned in this work to commit \releasecommit{} of the public release; no new density-functional calculation was performed. The target is the formation energy $E_\mathrm{f}$ of the neutral vacancy $V_\mathrm{O}^{\times}$ at the reference oxygen chemical potential of the release, computed with PBEsol(+$U$) in a supercell of the relaxed pristine host \citep{Kumagai2021insights}. Removing a neutral oxygen atom leaves two electrons behind, and whether they settle in a localised state on the neighbouring cations or in the host band edge is the electronic question that the descriptors encode.

The release covers \numNetReleaseOxides{} oxides and holds \numFilterCascadeStepsReleaseSites{} completed neutral calculations in \numFilterCascadeStepsReleaseHosts{} of them. Starting from the neutral machine-learning subset of Kumagai et al., three filters (Appendix~\ref{app:data}) remove the entries in which the oxygen did not remain a simple vacancy (split vacancies, unknown defect types and reconfigured final states) and one host whose site labels do not match the targets, and check the site mapping of two hosts with inconsistent cell metadata. The dataset that remains holds \numFilterCascadeStepsFinalSites{} sites in \numFilterCascadeStepsFinalHosts{} hosts. A host is one crystal structure with a distinct formula, and a site is one symmetry-inequivalent oxygen of it.

The release provides \numDescriptorProvenanceNDescriptorColumns{} descriptors per site, which we class by the calculation that a new oxide would need (Appendix~\ref{app:data}). The host-electronic class, \numNetDescriptorsHost{} host-level descriptors, holds the PBEsol(+$U$) band gap, the oxide formation energy per atom, the electronic and ionic dielectric constants and the $s$, $p$, $d$ and $f$ character of the valence-band maximum and conduction-band minimum. The site-electronic class, \numNetDescriptorsSite{} site-level descriptors, holds the Bader charge and volume of the oxygen to be removed \citep{Henkelman2006fast}, its Born effective charge, the corresponding statistics over its cation neighbours and the centre of its oxygen $p$ band. These \numRfElectronicNFeatures{} pristine-cell electronic descriptors, \numNetDescriptorsHost{} host-level and \numNetDescriptorsSite{} site-level, are given to D. The \numRfStructuralNFeatures{} structural and compositional descriptors, the cation electronegativities and the composition, are functions of the structure that every network already reads; they are withheld from D, so that S and D differ only in electronic information. Every descriptor of D comes from the pristine cell and none requires a defect calculation, so D is an oracle for the electronic structure of the pristine host, not for the defect.

The perturbed-host-state filter of the release sets the domain of every result. In a perturbed-host-state entry the two electrons of the neutral vacancy do not occupy a localised gap state but perturb the host band edge, so the vacancy is a shallow donor whose electrons delocalise into the conduction band \citep{Lany2005anion,Zhang2006systematic}; a supercell energy of such a configuration does not describe an isolated neutral defect, and Kumagai et al.\ excluded these entries from their machine-learning subset \citep{Kumagai2021insights}. Of the \numScreeningAbsenceReleaseNeutralCompletedEntries{} completed neutral entries, \numScreeningAbsenceReleaseNeutralCompletedPhs{} carry the flag, and they concentrate on \dzero{} cations with a low-lying empty $d$ band: the flagged fraction is \numScreeningAbsencePhsByCationNbPhsRateWith{} for hosts containing Nb, \numScreeningAbsencePhsByCationWPhsRateWith{} for W and \numScreeningAbsencePhsByCationTiPhsRateWith{} for Ti, against \numScreeningAbsencePhsByCationZrPhsRateWith{} for Zr, whose empty $d$ band lies higher, and \numScreeningAbsencePhsByCationGaPhsRateWith{} for Ga. SrTiO$_3$ and BaTiO$_3$ leave the dataset for this reason. Every result of this paper therefore describes deep, localised neutral oxygen vacancies in ordered oxides classified as non-magnetic in the release, computed with one functional and one workflow. Shallow-donor oxides lie outside the domain of every model, and charged vacancies were not tested.
